\subsection{A conformally invariant coordinate test on a selected collar}
\label{sec:disk-collar-comparison}

Take a marked bicollar of one of the actual disks in
Theorem~\ref{thm:endpoint-disks} or Theorem~\ref{thm:band-charge}, contained in its open model
ball. Use the disk's angular parametrization and send its
boundary to the unit sphere of \(A=\{e^{-1}<|x|<e\}\) by an
increasing radial change of parameter. For the resulting
embedding \(\Psi:A\hookrightarrow M\), put
\[
 m=(\det p)^{-1/4}p=\exp h,\qquad p=\Psi^*k_i,\qquad
 G_\alpha(m)=\partial_j(m^{\alpha\alpha}m^{j\alpha})
 \quad\text{(no sum on \(\alpha\))}.
\]
The same \(m\) arises from the physical pullback \(\Psi^*g_i\).
Before normalization the coordinate flux has the expression
\(\sqrt{\det p}\,p^{\alpha\alpha}p^{j\alpha}\), which is
unchanged by the conformal multiplier. Thus \(G=0\) is the
four-harmonic coordinate condition in either gauge. The
determinant-one and \(n\)-harmonic gauge is the one of
Liimatainen--Salo \cite[Theorem~1.1]{LiimatainenSalo2016ConformalGauge}.

\begin{proposition}[Physical Weyl bound and differentiated coordinate residual]
\label{prop:disk-collar-comparison}
Fix \(A_0\Subset A\) around the seam and a straightening cutoff
\(0\le\psi\le1\) supported in \(A_0\), equal to one near the
seam. Set \(a=\|h\|_{\infty,A}\), \(\tau=\|DG(m)\|_{2,A}\),
and \(w^2=4(E_+(g_i;\Psi(A))+E_-(g_i;\Psi(A)))\). Derivative
norms use the Euclidean components and measure on the indicated
annulus. For \(a\) within a fixed small threshold, write
\(b=w+\tau+a\). Then
\[
 \|D^2h\|_{2,A_0}\le Cb,\qquad
 \|Dh\|_{4,A_0}^2\le Cab .
\]
With the source's scalar multiplier restored, the path
\(n_s=\exp((1-s\psi)h)\) gives the joined and exterior
splice estimates of Theorem~\ref{thm:simultaneous} using the
per-collar error
\begin{equation}\label{eq:disk-collar-allowance}
 \beta=C(a+\sqrt{ab}+ab).
\end{equation}
The constants depend on the fixed trim and profiles.
\end{proposition}

\begin{proof}
The matrix \(h\) is trace-free. At the Euclidean metric let
\(\mathcal W_0\) denote linearized covariant Weyl and put
\(\mathcal G_{0,\alpha}h=-\partial_jh_{j\alpha}
-\partial_\alpha h_{\alpha\alpha}\), with no sum on \(\alpha\).
At \(\xi\ne0\), the Weyl symbol kernel consists of
\(\ell\Eucl+\xi\otimes u+u\otimes\xi\). Indeed, in a frame with
first direction parallel to \(\xi\), the Riemann symbol has only
the components generated by \(R'_{1a1b}\), \(a,b>1\). If its
Weyl part vanishes, it is the Kulkarni--Nomizu product of the
Euclidean metric with a Schouten tensor. Restriction to the
three-dimensional transverse space forces the transverse
Schouten block to vanish. The mixed components force its
\(1a\) entries to vanish, leaving \(R'_{1a1b}\) proportional to
\(\delta_{ab}\). The converse follows from conformal and
diffeomorphism invariance. For the trace-free kernel,
\(4\ell+2\xi\cdot u=0\), and
\[
 \xi_j h_{j\alpha}+\xi_\alpha h_{\alpha\alpha}
 =(|\xi|^2+2\xi_\alpha^2)u_\alpha .
\]
Thus \((\mathcal W_0,D\mathcal G_0)\) is jointly elliptic on
trace-free tensors. Plancherel bounds \(\|D^2f\|_2\) by the
\(L^2\) norms of this pair for compactly supported trace-free \(f\).

Use this bound on \(\chi^2h\), where \(\chi\) is a fixed
interior cutoff equal to one on \(A_0\), and set
\(X=\|\chi^2D^2h\|_2\) and \(J=\|\chi Dh\|_4^2\).
The nonlinear errors in covariant Weyl and \(DG\) are bounded
by \(C(a|D^2h|+|Dh|^2)\), by smooth dependence of their
principal coefficients on \(h\). The covariant Weyl norm in
coordinates is at most \(Cw\) by conformal invariance and
uniform ellipticity. The first-order cutoff terms carry one
factor of \(\chi\). It follows that
\[
 X\le C(b+\|\chi Dh\|_2+aX+J).
\]
Integration by parts in \(\int\chi^2|Dh|^2\) gives
\(\|\chi Dh\|_2\le C\sqrt{aX}+Ca\). Applying the same method
to \(\int\chi^4|Dh|^4\) gives
\(J^2\le Ca(JX+J^{3/2})\), hence \(J\le Ca(X+a)\).
Young's inequality and small fixed \(a\) give the two
interior bounds.

In the Euclidean density representation of chiral Weyl
used in Theorem~\ref{thm:simultaneous}, the matrix
\(F_\pm(\exp f)\) is a constant-coefficient linear function
of \(D^2f\) up to \(O(|f|\,|D^2f|+|Df|^2)\). In the
difference \(F_\pm(n_s)-(1-s\psi)F_\pm(m)\) the linear
second derivatives cancel except for cutoff terms of order
at most one in \(h\). The remaining terms are bounded in
\(L^2\) by \(C(a\|D^2h\|_2+\|Dh\|_4^2+\|Dh\|_4+a)\) on
the fixed support. This gives \eqref{eq:disk-collar-allowance}
on the whole support and its exterior half. Restoring the
scalar multiplier preserves the chiral energies. At \(s=1\)
the metric is Euclidean in comparison gauge near the seam,
so the same round-product insertion and the same original
disk marking give the stretched and capped comparisons.
\end{proof}

For disjoint selected collars, \(w_j^2<\zeta\) and
\(\sum_jw_j^2\le8e_i\). Put \(A_2=(\sum_j a_j^2)^{1/2}\)
and \(J_1=\sum_j a_j(w_j+\tau_j+a_j)\). Taking the direct
sum in \eqref{eq:disk-collar-allowance} gives
\(\eta\le C(A_2+\sqrt{J_1}+J_1)\). The physical contribution
\(\sum_j a_jw_j\) is bounded by \(A_2\sqrt{8e_i}\).
In particular \(A_2\to0\) and \(\sum_j a_j\tau_j\to0\)
are sufficient for the joined energy bound \(e_i+o(1)\)
and the retained bound \(e_i-\sum_j E_\pm(g_i;B_{ij})+o(1)\).
The first bound keeps the joined path strictly below \(2T\),
so the capped class is positive. For \(I<4T/3\) choose a tolerance smaller than
\((4T-3I)/2\) in Theorem~\ref{thm:band-charge}.
A fixed finite family then carries a uniformly positive amount
of either specified chirality. By passing to a subsequence some
fixed disk index has charge bounded away from zero for that sign.
A vanishing allowance on this charged disk alone gives a positive
same-type competitor with energy below \(I\) at large indices.
More generally it suffices that the measured savings exceed
\(e_i-I\) at one index and that the joined bound is strictly
below \(2T\). The global equality of the two chiral energies
allows the same argument using either sign.

The disk sandwich does not control these coordinate tests
for arbitrary parametrizations. For example, in a round
stereographic chart compose with the increasing radial map
\(r\mapsto r+N^{-3/2}\chi_1(r)\sin(N(r-1))\), with fixed
\(\chi_1=1\) near one and supported between both collar
faces. It fixes the seam and both faces, preserving every
source ball inclusion and curvature integral. If \(f_N\)
is this radial map and \(\rho_N=f_N'/(f_N/r)\), the
determinant-one metric has radial eigenvalue
\(\rho_N^{3/2}\) and tangential eigenvalue \(\rho_N^{-1/2}\).
The distortion tends to zero as \(N^{-1/2}\), but
\(\|Dh_N\|_{4,A_0}\ge cN^{1/2}\) from the radial derivative
of \(\log\rho_N\). Since Weyl vanishes, the interior estimate
implies \(a_N\tau_N\to\infty\). The ordinary round chart
straightens with zero cost. This example concerns the
prescribed coordinates. Finding a controlled annulus
retaining an actual charged disk remains an additional
geometric requirement.
